\documentclass[11pt]{article}
\usepackage{fontspec}
\setmainfont{Times New Roman}
\setsansfont{Arial}
\setmonofont{Consolas}
\usepackage{amsmath,amssymb}
\usepackage{geometry}
\usepackage{microtype}
\geometry{a4paper,margin=3cm}
\setlength{\parindent}{1.25em}
\setlength{\parskip}{0pt}
\emergencystretch=18em
\allowdisplaybreaks
\raggedbottom
\providecommand{\vdotswithin}[1]{\mathrel{\vdots}}

\begin{document}

\section*{\S\ 6. Najobčejše oblasti iz cělyh transcendentnyh čisel pri osnovanju na racionalnoj bazi.}

V analogiji s algebraičnoj bazoj racionalnu bazu trěba definirati tako:
«Sistem $\Theta$ čisel $\vartheta$ imenuje se racionalna baza vsih čisel, ako
$\Theta$ dozvoljaje izraziti vsako čislo racionalno --- s koeficientami iz
$K$\footnote{Pod racionalnymi funkcijami vždy razumějut se take s algebraičnymi
čislovymi koeficientami; te algebraične čislove koeficienty vključene sut do
definicije zaradi III i IV. Bole odgovarjajuče bylo by trěbovati racionalne
čislove koeficienty kako v III i IV, tako i pri definiciji racionalnoj bazy.
Zaključenja iz \S\ 6 i \S\ 7 pri tom ostavajut doslovno ty same, ako samo polje
$K$ vsih algebraičnyh čisel zaměnimo poljem vsih racionalnyh čisel.} --- togda
kak pri najmanje jednom dobrom uporjadkovanju nijedno bazovo čislo ne može se
izraziti racionalno --- s koeficientami iz $K$ --- črez konečno mnogo
prědhodnyh bazovyh čisel».\footnote{V protivnosti k linearnoj i algebraičnoj
bazi, tu porjadok elementov igraje rolju. V porjadku
$\eta,\sqrt[\nu]{\eta}$, napr., trěba vzęti oba elementy, a v obratnom porjadku
samo $\sqrt[\nu]{\eta}$.}

Nehaj oblast $\mathfrak G$ iz cělyh transcendentnyh čisel bude podložena
dobromu uporjadkovanju $\Omega$. Togda može se slučiti, že veličina $z$ iz
$\mathfrak G$ jest racionalno izrazima črez konečno mnogo prědhodnyh veličin iz
$\mathfrak G$, t. j. udovletvorjaje uravnenju:
$z=\varphi(\xi)$, gdě $\varphi(\xi)$ jest racionalna funkcija, s koeficientami
iz $K$, od $\xi_1,\ldots,\xi_t$; osoblivo vsako algebraično čislo $\alpha$
oblasti udovletvorjaje uravnenju $z=\alpha$. Ostale veličiny $\vartheta$ iz
$\mathfrak G$ tvoret sistem $\Theta$, ktory trěba dokazati kako racionalnu bazu
vsih veličin iz $\mathfrak G$. Bo po prědpoloženju vsaka veličina $\vartheta$
jest racionalno nezavisima od prědhodnyh veličin iz $\Theta$. Indukcijno
zaključenje\footnote{Srav. Zermelo, tamže, \S\ 1.} pokazuje dalej, že vsaka
relacija $z=\varphi(\xi)$ vleče za soboju relaciju $z=\psi(\vartheta)$; inače
moralo by byti prvo element
$z_0=\varphi_0(\xi_1,\ldots,\xi_t)$, ktory ne može se izraziti racionalno črez
$\vartheta$, togda kak prědhodne $\xi_1,\ldots,\xi_t$ sut racionalne funkcije od
$\vartheta$; iz togo proizhodi proturěčenje. $\Theta$ po II jest jednočasno
racionalna baza vsih čisel; po III i I $\mathfrak G$ mimo $\Theta$ soderži
oblast cělosti $[\Theta]$, sostaveču se iz vsih celočislenyh polinomov od
$\vartheta$, kako dělitelj, togda kak $[\Theta]$ po IV ne soderži nikakogo
algebraično-drobnogo čisla. Racionalna baza vsih čisel $\Theta$ zato
udovletvorjaje dodatnomu uslovju, že iz $\Theta$ izvedena oblast cělosti
$[\Theta]$ ne soderži nikakogo algebraično-drobnogo čisla;\footnote{Že to
faktično jest uslovje, vidi se napr. iz togo, že dvě veličiny
$\eta,\frac1{2\sqrt{\eta}}$ ne sut dopustime kako bazove elementy, no veličiny
$\eta,\frac1{\sqrt{\eta}}$ sut dopustime.} $\Theta$ jest {\itshape racionalna
baza s dodatnym uslovjem} za vsa čisla. Zato v analogiji s \S\ 2 važi:
{\itshape Vsaka oblast $\mathfrak G$ može se konstruovati pomoću racionalnoj
bazy s dodatnym uslovjem tako, že $\mathfrak G$ soderži oblast cělosti
$[\Theta]$ kako dělitelj.}

Nehaj $\Theta$ bude koja-koli racionalna baza s dodatnym uslovjem za vsa čisla
--- jejino postojanje jest zato osigurjeno postojanjem oblastij
$\mathfrak G$. Nehaj $\mathfrak N(\mathfrak G)$ označa sovkupnost vsih oblastij
$\mathfrak G$, ktore soderžet $\Theta$ i poslědovateljno $[\Theta]$. Kogda
$\Theta$ proběgaje vsaku proizvoljnu bazu s dodatnym uslovjem, togda
$\mathfrak N(\mathfrak G)$ po vyše skazanom proběgaje sovkupnost vsih oblastij
$\mathfrak G$; zato možemo znovu ograničiti se na oblasti $\mathfrak G$ iz
$\mathfrak N(\mathfrak G)$.

Nyně samo $[\Theta]$ jest oblast $\mathfrak G$; bo vsako čislo može se izraziti
racionalno črez $\Theta$, znači kako količnik dvuh polinomov iz $[\Theta]$; i
$[\Theta]$ po svojoj konstrukciji udovletvorjaje takože ostalym uslovjam I, III,
IV. {\itshape $[\Theta]$ jest zato največši obči dělitelj ili prěsěk vsih
oblastij $\mathfrak G$ iz $\mathfrak N(\mathfrak G)$; dalje
$\mathfrak N(\mathfrak G)$ jest zatvorjena vzhodno k tvorenju prěsěkov, t. j.
takože prěsěk vsih oblastij $\mathfrak G$ vsakogo podsistema iz
$\mathfrak N(\mathfrak G)$ naleži k $\mathfrak N(\mathfrak G)$.} Bo za prěsěk
jest ispolnjeno I; II i III slědujut iz togo, že $[\Theta]$ jest dělitelj; IV iz
togo, že prěsěčna oblast soderži se v někoj oblasti $\mathfrak G$ kako
dělitelj.

Racionalno-cěla adjunkcija proizvoljnogo sistema $\mathfrak S$ racionalnyh
funkcij od $\vartheta$ k $[\Theta]$ --- vsaka algebraična funkcija od
$\vartheta$ može se, pomoću svojstva racionalnoj bazy, racionalno izraziti črez
někotore druge $\vartheta$ --- daje oblast cělosti $[\Theta,\mathfrak S]$,
ktora imaje svojstva I, II, III i zato stavaje oblast $\mathfrak G$, skoro
$\mathfrak S$ jest «dopustimy sistem», t. j. udovletvorjaje uslovju, že črez
racionalno-cělu adjunkciju nikako algebraično-drobno čislo ne vhodi v oblast.
Obratno, vsaka oblast $\mathfrak G$ iz $\mathfrak N(\mathfrak G)$ dopustjaje
predstavjenje $[\Theta,\mathfrak S]$, ako samo $\mathfrak S$ označa ostatny
sistem $\mathfrak G-[\Theta]$. Zato za najobčejše oblasti $\mathfrak G$ v polnoj
analogiji s \S\ 4 važe rezultaty:
\begin{enumerate}
\item[a)] {\itshape Prěsěk vsih oblastij $\mathfrak G$ iz
$\mathfrak N(\mathfrak G)$ jest dan oblastju cělosti $[\Theta]$, ktora sama
jest oblast $\mathfrak G$; $\mathfrak N(\mathfrak G)$ jest zatvorjena vzhodno k
tvorenju prěsěkov.}
\item[b)] {\itshape Vsaka oblast $\mathfrak G$ iz $\mathfrak N(\mathfrak G)$
nastavaje črez racionalno-cělu adjunkciju dopustimogo sistema $\mathfrak S$ k
$[\Theta]$. Kogda $\Theta$ proběgaje vsaku možnu racionalnu bazu s dodatnym
uslovjem, $\mathfrak N(\mathfrak G)$ proběgaje sovkupnost vsih oblastij
$\mathfrak G$.}\footnote{O najobčejših dopustimyh sistemah $\mathfrak S$ važi
to, čto bylo skazano v \S\ 4, ako samo «algebraično» zaměnimo slovom
«racionalno».}
\end{enumerate}

{\itshape Primětka:} Ako iz $\Theta$ izberemo algebraičnu bazu $H$ vsih veličin
iz $\Theta$, togda $H$ jednočasno stavaje algebraična baza vsih čisel; obratno,
vsaku algebraičnu bazu možno doplniti do racionalnoj, kladuči $H$ na početok
dobrogo uporjadkovanja. Poslě togo konstrukciju oblastij $\mathfrak G$, danu v
\S\ 5 pomoću algebraičnoj bazy, možno razuměti tako: sistem $\mathfrak S$, ktory
trěba adjungovati k $[H]$, razdělja se na dva podsistemy $\mathfrak S_1$ i
$\mathfrak S_2$; adjunkcija $\mathfrak S_1$ svodi se k dopolnjenju $H$ do
racionalnoj bazy s dodatnym uslovjem, k kotoroj potom racionalno-cělo
adjunguje se dopustimy sistem $\mathfrak S_2$.

\end{document}
